% ============================================================
% CHAPITRE 5 — Sprint 2 : Organisation Admin
% ============================================================
\chapter{Sprint 2 --- Organisation Admin (Zones, Agences et Affectation RH)}

\section{Introduction}

Après la mise en place de l'identité (Sprint~1), le Sprint~2 structure
l'organisation territoriale de la plateforme \textbf{DAAM}. Le recrutement
est rattaché à des \textbf{zones} et à des \textbf{agences} (entreprises) ;
un RH ne peut piloter des offres que dans les zones qui lui sont affectées.

Ce sprint se concentre donc sur trois briques Admin~: création des zones,
création des agences rattachées à une zone, et affectation exclusive d'un
RH à une ou plusieurs zones. Ces fondations conditionnent ensuite la
création des offres (Sprint~3).


\section{Objectifs du Sprint}

Les principaux objectifs de ce sprint sont~:
\begin{itemize}
  \item permettre à l'administrateur de créer et lister les zones
        géographiques / organisationnelles~;
  \item permettre de créer des agences (nom, adresse, localisation) liées
        à une zone~;
  \item affecter un compte RH à une ou plusieurs zones, avec la règle
        \textbf{un seul RH par zone}~;
  \item exposer ces données via le microservice Recruitment et les
        interfaces Admin correspondantes~;
  \item garantir que, dès ce sprint, le périmètre RH est borné par ses
        affectations (filtrage des agences visibles).
\end{itemize}


\section{Spécification et Analyse des Besoins}

Cette section détaille les exigences et scénarios qui ont guidé le
développement. Le Sprint Backlog décompose les User Stories en tâches,
tandis que le diagramme de cas d'utilisation illustre les interactions
Admin (et la consultation RH limitée).


\subsection{Sprint Backlog}

\begin{table}[htbp]
  \centering
  \scriptsize
  \caption{Sprint Backlog du Sprint 2}
  \label{tab:sprint2-backlog}
  \begin{tabular}{|p{2.2cm}|c|p{4.0cm}|c|p{3.3cm}|c|}
    \hline
    \textbf{Fonctionnalité} & \textbf{ID} & \textbf{Scénario utilisateur} &
    \textbf{ID Tâche} & \textbf{Tâches} & \textbf{Est.} \\
    \hline
    Zones
    & US-04
    & En tant qu'Admin, je crée et consulte les zones.
    & T1
    & API \texttt{POST/GET /zones}
    & 1,5 j \\
    \cline{4-6}
    & & & T2 & Interface \texttt{/admin/zones} & 1 j \\
    \hline
    Agences
    & US-05
    & En tant qu'Admin, je crée une agence rattachée à une zone (adresse, carte).
    & T3
    & API \texttt{POST/GET /companies}
    & 2 j \\
    \cline{4-6}
    & & & T4 & Interface \texttt{/admin/companies} & 1,5 j \\
    \hline
    Affectation RH
    & US-06
    & En tant qu'Admin, j'affecte un RH à une ou plusieurs zones (1 RH / zone).
    & T5
    & API \texttt{POST/GET /rh-zone-assignments}
    & 2 j \\
    \cline{4-6}
    & & & T6 & Interface \texttt{/admin/rh-assignments} & 1,5 j \\
    \hline
    Périmètre RH
    & US-07
    & En tant que RH, je ne vois que les agences de mes zones affectées.
    & T7
    & Filtrage \texttt{GET /companies} selon affectations
    & 1 j \\
    \hline
  \end{tabular}
\end{table}


\subsection{Diagramme de cas d'utilisation du Sprint}

\begin{figure}[H]
  \centering
  \reportfig{img/uc-sprint2.png}
  \caption{Diagramme de cas d'utilisation du Sprint 2}
  \label{fig:uc-sprint2}
\end{figure}


\section{Description textuelle des cas d'utilisation}

\subsection{Description textuelle du cas d'utilisation <<~Ajouter une zone~>>}

\begin{table}[htbp]
  \centering
  \small
  \caption{Description textuelle du cas d'utilisation <<~Ajouter une zone~>>}
  \label{tab:uc-add-zone}
  \begin{tabular}{|p{3.0cm}|p{10.2cm}|}
    \hline
    \textbf{Titre} & Ajouter une zone \\
    \hline
    \textbf{Acteur} & Administrateur \\
    \hline
    \textbf{Précondition} & L'Admin est authentifié (\texttt{ROLE\_ADMIN}). \\
    \hline
    \textbf{Postcondition} & Une zone active est créée (identifiant \texttt{zoneId}). \\
    \hline
    \textbf{Scénario nominal}
    &
    \begin{enumerate}[leftmargin=*,itemsep=0.15em,topsep=0.2em]
      \item L'Admin ouvre \texttt{/admin/zones}.
      \item Il saisit le nom (unique) et éventuellement une description.
      \item Le frontend envoie \texttt{POST /api/recruitment/zones}.
      \item Le microservice Recruitment enregistre la zone.
      \item La liste des zones est mise à jour.
    \end{enumerate}
    \\
    \hline
    \textbf{Alternatives}
    & Nom déjà utilisé~: création refusée. \\
    \hline
  \end{tabular}
\end{table}


\subsection{Description textuelle du cas d'utilisation <<~Ajouter une agence~>>}

\begin{table}[htbp]
  \centering
  \small
  \caption{Description textuelle du cas d'utilisation <<~Ajouter une agence~>>}
  \label{tab:uc-add-company}
  \begin{tabular}{|p{3.0cm}|p{10.2cm}|}
    \hline
    \textbf{Titre} & Ajouter une agence \\
    \hline
    \textbf{Acteur} & Administrateur \\
    \hline
    \textbf{Précondition} & Au moins une zone existe. \\
    \hline
    \textbf{Postcondition} & Une agence est créée et rattachée à une zone. \\
    \hline
    \textbf{Scénario nominal}
    &
    \begin{enumerate}[leftmargin=*,itemsep=0.15em,topsep=0.2em]
      \item L'Admin ouvre \texttt{/admin/companies}.
      \item Il choisit la zone, saisit le nom, l'adresse et éventuellement
            les coordonnées / lien Google Maps.
      \item Le frontend envoie \texttt{POST /api/recruitment/companies}.
      \item Le service persiste l'agence avec \texttt{companyId} et \texttt{zoneId}.
      \item L'agence apparaît dans la liste.
    \end{enumerate}
    \\
    \hline
  \end{tabular}
\end{table}


\subsection{Description textuelle du cas d'utilisation <<~Affecter un RH a des zones~>>}

\begin{table}[htbp]
  \centering
  \small
  \caption{Description textuelle du cas d'utilisation <<~Affecter un RH a des zones~>>}
  \label{tab:uc-assign-rh}
  \begin{tabular}{|p{3.0cm}|p{10.2cm}|}
    \hline
    \textbf{Titre} & Affecter un RH a des zones \\
    \hline
    \textbf{Acteur} & Administrateur \\
    \hline
    \textbf{Précondition} & Un compte \texttt{ROLE\_RH} existe~; des zones sont disponibles. \\
    \hline
    \textbf{Postcondition} & Le RH est lié aux zones choisies. \\
    \hline
    \textbf{Scénario nominal}
    &
    \begin{enumerate}[leftmargin=*,itemsep=0.15em,topsep=0.2em]
      \item L'Admin ouvre \texttt{/admin/rh-assignments}.
      \item Il sélectionne un RH et une ou plusieurs zones encore libres.
      \item Le frontend envoie \texttt{POST /api/recruitment/rh-zone-assignments}.
      \item Le service vérifie que le compte est bien RH et qu'aucune
            autre affectation n'existe déjà pour chaque zone.
      \item Les affectations sont enregistrées.
    \end{enumerate}
    \\
    \hline
    \textbf{Alternatives}
    & Zone déjà attribuée à un autre RH, ou utilisateur non RH~: erreur. \\
    \hline
  \end{tabular}
\end{table}


\subsection{Autres cas d'utilisation}

\begin{itemize}
  \item \textbf{Lister les zones / agences / affectations}~: l'Admin
        consulte l'organisation complète.
  \item \textbf{Consulter les agences de son périmètre}~: le RH appelle
        \texttt{GET /companies} et ne reçoit que les agences des zones
        qui lui sont affectées (préparation du Sprint~3).
\end{itemize}


\section{Conception}

\subsection{Diagramme de classes}

Le domaine du sprint~2 s'articule autour de \textbf{Zone}, \textbf{Company}
(agence) et \textbf{RhZoneAssignment}. Une agence appartient à une zone ;
une affectation relie un \texttt{rhUserId} à un \texttt{zoneId}.

\begin{figure}[H]
  \centering
  \reportfig{img/class-sprint2.png}
  \caption{Diagramme de classes du Sprint 2}
  \label{fig:class-sprint2}
\end{figure}


\subsection{Diagrammes de séquence}

\subsubsection{Diagramme de séquence <<~Ajouter une zone~>>}

\begin{figure}[H]
  \centering
  \reportfig{img/seq-add-zone.png}
  \caption{Diagramme de s\'equence --- Ajouter une zone}
  \label{fig:seq-add-zone}
\end{figure}

\subsubsection{Diagramme de séquence <<~Ajouter une agence~>>}

\begin{figure}[H]
  \centering
  \reportfig{img/seq-add-company.png}
  \caption{Diagramme de s\'equence --- Ajouter une agence}
  \label{fig:seq-add-company}
\end{figure}

\subsubsection{Diagramme de séquence <<~Affecter un RH a des zones~>>}

\begin{figure}[H]
  \centering
  \reportfig{img/seq-assign-rh.png}
  \caption{Diagramme de s\'equence --- Affecter un RH a des zones}
  \label{fig:seq-assign-rh}
\end{figure}


\section{Réalisation}

\subsection{Interface de gestion des zones}

L'espace Admin permet de créer une zone (nom, description) et d'en
consulter la liste.

\begin{figure}[H]
  \centering
  \reportfig{img/ui-zones.png}
  \caption{Interface de gestion des zones (Admin)}
  \label{fig:ui-zones}
\end{figure}


\subsection{Interface de gestion des agences}

La page agences permet de rattacher une entreprise à une zone et de
saisir l'adresse / le lien cartographique.

\begin{figure}[H]
  \centering
  \reportfig{img/ui-companies-1.png}
  \caption{Interface de gestion des agences (Admin) --- liste}
  \label{fig:ui-companies}
\end{figure}

\begin{figure}[H]
  \centering
  \reportfig{img/ui-companies-2.png}
  \caption{Interface de gestion des agences (Admin) --- formulaire}
  \label{fig:ui-companies-form}
\end{figure}


\subsection{Interface d'affectation RH}

L'Admin sélectionne un utilisateur RH et les zones encore libres. La
règle métier impose un unique RH par zone.

\begin{figure}[H]
  \centering
  \reportfig{img/ui-rh-assign.png}
  \caption{Interface d'affectation RH -- zones (Admin)}
  \label{fig:ui-rh-assign}
\end{figure}


\section{Conclusion}

Le Sprint~2 a posé l'organisation territoriale de DAAM~: zones, agences
et affectations RH. L'administrateur pilote la structure~; le RH est
ensuite limité à son périmètre.

Le chapitre suivant présente le Sprint~3, consacré aux
\textbf{offres de recrutement} et aux \textbf{candidatures}.

\clearpage